\section{Konsep Usulan Riset Lanjutan: Physics-Constrained Spectral Transformer}

\subsection{Judul Konseptual}
Pengembangan \textit{Physics-Constrained Spectral Transformer} (PCST): \textit{Framework Physics-Informed Deep Learning} untuk analisis spektrum LIBS dengan pendekatan \textit{calibration-free}.

\subsection{Ringkasan Konsep}
Konsep riset lanjutan ini mengusulkan \textit{Physics-Constrained Spectral Transformer} (PCST) sebagai upaya mitigasi masalah tumpang tindih spektrum LIBS pada sampel kompleks. Berbeda dari pendekatan \textit{convolutional neural network} konvensional, PCST dirancang untuk memanfaatkan arsitektur transformer dalam menangkap korelasi spektral jarak jauh (\textit{long-range dependencies}). Untuk menangani ribuan kanal spektrum secara efisien, mekanisme perhatian jarang (\textit{sparse attention}) berbasis keluarga Informer diposisikan sebagai kandidat utama. Model ini diusulkan untuk mengintegrasikan simulator transfer radiatif LTE/\textit{partial}-LTE sebagai kendala fisika, sehingga proses rekonstruksi profil emisi murni dan kuantifikasi unsur tetap berada pada manifold yang dipandu hukum fisika. Luaran yang ditargetkan mencakup prototipe perangkat lunak, dataset sintetik terstruktur, validasi laboratorium, dan publikasi ilmiah.

\subsection{Urgensi dan Permasalahan}
Laser-Induced Breakdown Spectroscopy (LIBS) menawarkan analisis unsur yang cepat, \textit{in-situ}, dan minim preparasi, sehingga relevan untuk inspeksi industri, pemantauan material, dan karakterisasi sampel geologi. Namun, plasma LIBS bersifat transien, tidak homogen, dan rentan menghasilkan \textit{severe spectral overlap}, terutama pada sampel alami yang kaya unsur. Distorsi ini diperparah oleh pelebaran Stark dan Doppler serta efek \textit{self-absorption} pada kondisi optis tebal.

Dalam konteks tersebut, terdapat tiga masalah pokok. Pertama, terdapat dilema akurasi dan kecepatan: inversi spektrum padat dengan fitting fisika penuh sangat mahal secara komputasi. Kedua, model AI \textit{black box} sering tidak mematuhi kaidah fisika plasma dan mudah gagal di domain data baru. Ketiga, \textit{self-absorption} masih menjadi penghalang utama bagi kuantifikasi LIBS bebas kalibrasi karena merusak linearitas hubungan antara intensitas garis dan konsentrasi unsur. Oleh karena itu, diperlukan arsitektur yang mampu mempertahankan efisiensi inferensi sekaligus tetap tunduk pada kendala transfer radiatif.

\subsection{Pendekatan Pemecahan Masalah}
Pendekatan yang diusulkan adalah \textit{physics-informed machine learning} berbasis transformer spektral. Secara konseptual, strategi pemecahan masalah meliputi tiga komponen. Pertama, penggunaan transformer spektral untuk menangkap konteks global dari seluruh rentang panjang gelombang, menggantikan keterbatasan reseptif lokal dari CNN. Kedua, integrasi kendala fisika radiatif ke dalam fungsi rugi melalui residu persamaan transfer radiatif, kendala opasitas, dan konsistensi termodinamika. Ketiga, penggunaan skema \textit{cycle consistency} atau rekonstruksi ulang spektrum sebagai validasi mandiri agar spektrum hasil prediksi tetap konsisten dengan spektrum masukan.

Dalam kerangka implementasi yang realistis, manuskrip utama pada repositori ini menyediakan fondasi forward model dua-zona \textit{partial}-LTE dengan transfer radiatif deterministik. Oleh sebab itu, PCST tidak diposisikan sebagai model yang berdiri sendiri, melainkan sebagai perluasan dari generator spektrum fisika yang telah dibangun. Dengan demikian, model transformer diarahkan untuk belajar di atas manifold sintetis yang telah dikendalikan secara fisis, bukan di atas data yang sepenuhnya tidak terstruktur.

\subsection{Bentuk Analitik Fondasi Fisika}
Fondasi analitik usulan ini tetap bertumpu pada model maju dua-zona yang terdiri atas zona inti (\textit{core}) dan zona selubung (\textit{shell}). Transport radiasi pada medium bergradien dinyatakan oleh persamaan transfer radiatif
\begin{equation}
\frac{dI_\lambda}{ds} = j_\lambda - \kappa_\lambda I_\lambda ,
\end{equation}
dengan $I_\lambda$ intensitas spektral, $j_\lambda$ koefisien emisi, dan $\kappa_\lambda$ koefisien absorpsi. Untuk konfigurasi dua-zona, bentuk solusi analitik yang dipakai pada generator sintetik adalah
\begin{equation}
I_{\mathrm{obs}}(\lambda)
= I_{\mathrm{core}}(\lambda)e^{-\tau_\lambda}
+ \frac{j_{\mathrm{shell}}(\lambda)}{\kappa_{\mathrm{shell}}(\lambda)}
\left(1-e^{-\tau_\lambda}\right),
\end{equation}
dengan
\begin{equation}
\tau_\lambda = \kappa_{\mathrm{shell}}(\lambda)\, d_{\mathrm{shell}} .
\end{equation}
Formulasi ini memungkinkan efek \textit{self-absorption} dan \textit{self-reversal} dimodelkan secara eksplisit tanpa fitting numerik berulang pada tahap inferensi.

Pada tingkat termodinamika, distribusi fraksi ionisasi dapat diturunkan dari hubungan Saha
\begin{equation}
\frac{n_e n_{II}}{n_I}
=
2\frac{Z_{II}(T_e)}{Z_I(T_e)}
\left(\frac{2\pi m_e k_B T_e}{h^2}\right)^{3/2}
\exp\left(-\frac{E_{\mathrm{ion}}}{k_B T_e}\right),
\end{equation}
sedangkan populasi tingkat eksitasi dalam pendekatan \textit{partial}-LTE dinyatakan melalui distribusi Boltzmann
\begin{equation}
n_i = \frac{n_{\mathrm{total}} g_i \exp(-E_i/k_B T_e)}{Z(T_e)}.
\end{equation}
Dengan demikian, generator dapat memetakan parameter plasma $\{T_e, n_e, \mathbf{c}, d_{\mathrm{core}}, d_{\mathrm{shell}}\}$ ke spektrum sintetis yang konsisten secara fisika.

Untuk pembentukan garis spektral, koefisien emisi dan absorpsi dituliskan sebagai
\begin{equation}
j_\lambda = \sum_{k\rightarrow i}\frac{hc}{4\pi\lambda}A_{ki}n_k\phi_V(\lambda),
\end{equation}
\begin{equation}
\kappa_\lambda =
\sum_{i\rightarrow k}
\frac{\lambda^4}{8\pi c}
\left(\frac{g_k}{g_i}\right)
A_{ki} n_i \phi_V(\lambda)
\left[1-\frac{n_k g_i}{n_i g_k}\right],
\end{equation}
dengan $\phi_V(\lambda)$ adalah profil Voigt hasil konvolusi pelebaran Doppler dan Stark. Pada level observasi instrumen, spektrum akhir masih dikonvolusikan dengan fungsi aparatus Gaussian
\begin{equation}
I_{\mathrm{final}}(\lambda)=I_{\mathrm{obs}}(\lambda)\otimes G_{\mathrm{inst}}(\lambda,\Delta\lambda_{\mathrm{inst}}).
\end{equation}

\subsection{Formulasi Analitik PCST}
Dengan fondasi forward model di atas, PCST diposisikan sebagai pemetaan
\begin{equation}
f_\theta:\mathbf{x}\in\mathbb{R}^{L}\mapsto
\left(\hat{T}_e,\hat{n}_e,\hat{\mathbf{c}},\hat{d}_{\mathrm{core}},\hat{d}_{\mathrm{shell}}\right),
\end{equation}
dengan $\mathbf{x}$ adalah spektrum masukan sepanjang $L$ kanal panjang gelombang, $\theta$ parameter model, dan $\hat{\mathbf{c}}$ adalah vektor komposisi unsur terprediksi.

Pada tahap berikutnya, keluaran model dapat diproyeksikan kembali ke domain spektral melalui operator simulator fisika $\mathcal{F}$:
\begin{equation}
\hat{\mathbf{x}} = \mathcal{F}\!\left(\hat{T}_e,\hat{n}_e,\hat{\mathbf{c}},\hat{d}_{\mathrm{core}},\hat{d}_{\mathrm{shell}}\right).
\end{equation}
Skema ini memungkinkan pembentukan \textit{cycle consistency} antara spektrum input dan spektrum hasil rekonstruksi.

Arsitektur loss total secara konseptual dapat ditulis sebagai
\begin{equation}
\mathcal{L}_{\mathrm{total}}
=
\alpha \mathcal{L}_{\mathrm{reg}}
+ \beta \mathcal{L}_{\mathrm{comp}}
+ \gamma \mathcal{L}_{\mathrm{phys}}
+ \delta \mathcal{L}_{\mathrm{cycle}},
\end{equation}
dengan:
\begin{enumerate}
    \item $\mathcal{L}_{\mathrm{reg}}$ adalah loss regresi untuk parameter termodinamika,
    \begin{equation}
    \mathcal{L}_{\mathrm{reg}}
    =
    \mathrm{MSE}(T_e,\hat{T}_e)
    +
    \mathrm{MSE}(n_e,\hat{n}_e),
    \end{equation}
    \item $\mathcal{L}_{\mathrm{comp}}$ adalah loss komposisi unsur,
    \begin{equation}
    \mathcal{L}_{\mathrm{comp}} = \mathrm{MSE}(\mathbf{c},\hat{\mathbf{c}}),
    \end{equation}
    dengan tambahan proyeksi simplex agar
    \begin{equation}
    \sum_{i=1}^{M}\hat{c}_i = 1, \qquad \hat{c}_i \ge 0,
    \end{equation}
    \item $\mathcal{L}_{\mathrm{phys}}$ adalah penalti konsistensi fisika, yang dapat memuat residu transfer radiatif atau penalti terhadap solusi yang menghasilkan kedalaman optik tak realistis,
    \begin{equation}
    \mathcal{L}_{\mathrm{phys}}
    =
    \left\|
    \frac{dI_\lambda}{ds} - j_\lambda + \kappa_\lambda I_\lambda
    \right\|_2^2,
    \end{equation}
    \item $\mathcal{L}_{\mathrm{cycle}}$ adalah loss rekonstruksi spektral,
    \begin{equation}
    \mathcal{L}_{\mathrm{cycle}} = \|\mathbf{x}-\hat{\mathbf{x}}\|_2^2 .
    \end{equation}
\end{enumerate}

Jika seleksi fitur mRMR(MIQ) digunakan sebelum PCST, maka pemilihan fitur dilakukan dengan skor
\begin{equation}
\mathrm{MIQ}(X_i)=
\frac{I(X_i;Y)}
{\frac{1}{|S|}\sum_{X_j\in S} I(X_i;X_j)+\varepsilon},
\end{equation}
dengan $I(X_i;Y)$ menyatakan \textit{mutual information} fitur ke-$i$ terhadap target, dan penyebutnya merepresentasikan redundansi rata-rata terhadap himpunan fitur terpilih $S$.

\subsection{Desain Data dan Eksperimen}
Secara praktis, usulan ini berdiri di atas dataset sintetis yang telah dibagi berdasarkan kelompok komposisi dan tier resolusi spektral. Dua kelompok komposisi didefinisikan sebagai:
\begin{itemize}
    \item kelompok A: unsur mayor,
    \item kelompok B: unsur mayor + unsur jejak.
\end{itemize}
Tier resolusi pragmatis didefinisikan sebagai:
\begin{itemize}
    \item rendah: 4000 titik,
    \item menengah: 8000 titik,
    \item tinggi: 16000 titik.
\end{itemize}
Kombinasi keduanya membentuk kode dataset
\begin{equation}
\{A\text{-}L, A\text{-}M, A\text{-}H, B\text{-}L, B\text{-}M, B\text{-}H\}.
\end{equation}
Untuk setiap kode dataset, evaluasi dapat dilakukan minimal pada dua pipeline:
\begin{enumerate}
    \item \textit{plain} physics-informed model,
    \item mRMR(MIQ) + physics-informed model.
\end{enumerate}
Desain ini memungkinkan analisis terkontrol terhadap pengaruh kompleksitas komposisi, resolusi spektral, dan reduksi fitur terhadap akurasi prediksi.

\subsection{Kebaruan Konseptual}
Kebaruan utama usulan ini terletak pada penggabungan kecepatan inferensi model AI dengan rigiditas kendala fisika radiatif. Berbeda dari simulator transfer radiatif murni yang mengandalkan \textit{iterative fitting} untuk setiap sampel, PCST diarahkan untuk melakukan inferensi instan setelah proses pelatihan. Sebaliknya, berbeda dari \textit{deep learning} murni yang hanya mengejar akurasi statistik, PCST diposisikan untuk memasukkan residu transfer radiatif dan konsistensi spektral ke dalam proses optimasi.

Dengan framing tersebut, novelty usulan ini mencakup:
\begin{enumerate}
    \item penggunaan model maju dua-zona sebagai generator data dan pembatas manifold fisika;
    \item pengembangan model inversi spektral yang tidak hanya \textit{data-driven}, tetapi juga \textit{physics-constrained};
    \item formulasi menuju transformer spektral untuk kuantifikasi LIBS pada kondisi optis tebal;
    \item validasi silang antara data sintetik dan data eksperimen nyata.
\end{enumerate}

\subsection{Roadmap Penelitian}
Roadmap penelitian dapat diringkas ke dalam empat fase. Fase pertama adalah era kemometrik klasik dan CF-LIBS berbasis asumsi plasma homogen. Fase kedua adalah peralihan ke \textit{deep learning} murni untuk klasifikasi dan regresi spektral. Fase ketiga, yang menjadi lingkup pengembangan saat ini, adalah pembangunan fondasi \textit{physics-informed learning} melalui forward model dua-zona, dataset sintetik, dan model inversi awal. Fase keempat adalah arah jangka menengah menuju PCST yang lebih efisien, terkompresi, dan siap digunakan pada sistem analisis LIBS yang lebih dekat ke kebutuhan lapangan.

\subsection{Keterkaitan dengan Pipeline yang Sudah Ada}
Pipeline yang sedang dibangun saat ini telah mencakup generator dataset sintetik, model inversi berbasis \textit{physics-informed learning}, seleksi fitur mRMR(MIQ), serta konfigurasi eksperimen berdasarkan kelompok komposisi dan tier resolusi spektral. Dengan demikian, konsep PCST dapat dipahami sebagai tahap evolusi berikutnya dari pipeline ini. Forward model yang ada tetap berfungsi sebagai mesin fisika utama, sementara model inversi yang saat ini berbasis jaringan saraf hierarkis dapat dijadikan baseline sebelum transisi ke arsitektur transformer spektral.

\subsection{Validasi Eksperimental}
Secara konseptual, model yang dikembangkan diarahkan untuk divalidasi menggunakan sampel eksperimen tanah vulkanik yang memiliki referensi XRF dan CF-LIBS. Validasi ini penting agar riset tidak berhenti pada domain simulasi, tetapi memiliki jembatan menuju data eksperimen nyata. Dalam kerangka tersebut, XRF diposisikan sebagai acuan utama komposisi unsur, sedangkan CF-LIBS dipakai sebagai baseline analitik berbasis domain LIBS klasik. Posisi ilmiah yang aman untuk tahap ini adalah validasi awal (\textit{initial validation}) dan \textit{proof-of-method}, bukan klaim universalitas penuh pada seluruh jenis matriks geologi.

\subsection{Metode Evaluasi}
Kinerja model dievaluasi menggunakan galat akar kuadrat rata-rata relatif dalam persen terhadap rentang target uji:
\begin{equation}
\mathrm{RMSE}_{\%}
=
\frac{
\sqrt{\frac{1}{N}\sum_{i=1}^{N}(y_i-\hat{y}_i)^2}
}
{\max(y)-\min(y)}
\times 100\%.
\end{equation}
Metrik ini diterapkan untuk temperatur elektron, densitas elektron, dan komposisi unsur. Dalam implementasi praktis pipeline saat ini, interpretasi performa dinyatakan secara heuristik sebagai: baik untuk RMSE $\le 10\%$, cukup untuk $10\% < \mathrm{RMSE} \le 20\%$, dan lemah untuk RMSE $>20\%$.

\subsection{Arah Integrasi ke Manuskrip Utama}
File konsep ini disusun sebagai naskah terpisah agar tidak langsung mengubah struktur manuskrip utama. Bagian-bagian di dalamnya dapat dipindahkan secara selektif ke naskah utama sesuai kebutuhan, misalnya untuk memperkuat bagian \textit{Introduction}, \textit{State of the Art}, \textit{Methodology}, atau \textit{Future Work}. Dengan pendekatan ini, repositori tetap memiliki pemisahan yang jelas antara manuskrip utama yang berfokus pada forward model deterministik dan draft konseptual yang menjelaskan arah pengembangan menuju \textit{Physics-Constrained Spectral Transformer}.
